\chapter{Antioxidants}

\section*{Definition and Overview}
Antioxidants are substances that protect cells from damage caused by \textbf{free radicals} -- unstable molecules with unpaired electrons that are produced as a normal by-product of using oxygen for energy, and also by smoking, alcohol, pollution, and ultraviolet light. When free radicals accumulate faster than the body can neutralise them, the resulting imbalance is called \textbf{oxidative stress}, which has been linked to ageing and to the development of many chronic diseases.

Important antioxidants include vitamins C and E, beta-carotene, the minerals selenium and zinc, and plant compounds such as flavonoids and polyphenols. The body also manufactures its own antioxidant defences, including glutathione and the enzymes superoxide dismutase, catalase, and glutathione peroxidase. Because oxidative stress has been implicated in heart disease, cancer, and dementia, antioxidant supplements are heavily marketed. The scientific evidence, however, shows that high-dose supplements do not prevent these diseases and can sometimes cause harm; food remains the safest and most reliable source.

\section*{Epidemiology}
Antioxidant supplements are among the most widely used dietary supplements worldwide; a substantial proportion of adults take them regularly, usually in the belief that they protect against heart disease, cancer, and the effects of ageing. Intake of antioxidant vitamins from food varies widely with diet -- people who eat few fruit and vegetables have lower intakes, and specific groups such as the elderly, smokers, and people on restricted diets are at greater risk of deficiency.

True deficiency syndromes are now uncommon in developed countries but still occur: vitamin C deficiency (scurvy) is seen in people with very restricted diets, such as heavy alcohol dependence or extreme elimination diets, and selenium and zinc deficiency occur in regions with poor soil or in malnourished populations. Conversely, excessive supplement use is common, and its harms -- including bleeding risk, toxicity, and drug interactions -- are frequently unrecognised.

\section*{Aetiology and Pathophysiology}
Free radicals are generated continuously as cells burn oxygen to produce energy, and their production increases with smoking, excess alcohol, intense inflammation, infection, and sun exposure. Because they carry an unpaired electron, they react aggressively with nearby molecules, damaging cell membranes, proteins, and DNA -- the molecular signature of oxidative stress. The body counters this with an integrated antioxidant system:

\begin{itemize}
    \item \textbf{Vitamin C:} water-soluble; works in body fluids, recycles vitamin E, and is required for collagen synthesis
    \item \textbf{Vitamin E:} fat-soluble; embeds in cell membranes and interrupts the chain reaction of membrane damage
    \item \textbf{Beta-carotene and vitamin A:} carotenoids with antioxidant and provitamin A activity
    \item \textbf{Selenium and zinc:} essential cofactors of antioxidant enzymes, including glutathione peroxidase
    \item \textbf{Flavonoids and polyphenols:} plant compounds with antioxidant and anti-inflammatory properties
    \item \textbf{Endogenous systems:} glutathione, superoxide dismutase, and catalase form the body's own first line of defence
\end{itemize}

\section*{Clinical Features}
Healthy people who eat a balanced diet generally have no symptoms related to antioxidants, because intake is adequate and excess is excreted or stored. Recognisable features arise mainly with deficiency or excess:

\begin{itemize}
    \item \textbf{Vitamin C deficiency (scurvy):} bleeding and swollen gums, easy bruising, poor wound healing, and fatigue
    \item \textbf{Vitamin E deficiency:} rare, but causes nerve damage with numbness, weakness, and poor coordination
    \item \textbf{Selenium excess:} brittle, ridged nails, hair loss, and a garlicky breath odour
    \item \textbf{Zinc deficiency:} poor wound healing, hair loss, skin rash, and impaired taste
    \item \textbf{No reliable symptoms of ``antioxidant lack'':} general fatigue or low energy is far more often due to anaemia, thyroid disease, depression, or poor sleep
    \item \textbf{Excess:} stomach upset and diarrhoea with high-dose vitamin C; bleeding risk with high-dose vitamin E
\end{itemize}

\section*{Differential Diagnosis}
Because antioxidant supplements are taken by people seeking general health benefits, the clinical questions usually relate to unexplained symptoms or to the safety of the supplements themselves:

\begin{itemize}
    \item \textbf{Anaemia, thyroid disease, or depression:} far more common causes of fatigue than any antioxidant deficiency
    \item \textbf{A genuine vitamin deficiency} versus an \textbf{unbalanced or restricted diet}
    \item \textbf{Harm from high-dose supplements:} bleeding, nausea, or toxicity mistaken for a new illness
    \item \textbf{Drug interactions:} supplement effects mistaken for the side effects of prescribed medicines
    \item \textbf{Underlying disease:} symptoms such as easy bruising may reflect liver disease, blood disorders, or anticoagulant use rather than antioxidant status
\end{itemize}

\section*{When to Seek Medical Care}
See a doctor if you have symptoms that might suggest a nutritional deficiency, such as easy bruising, bleeding gums, unexplained tiredness, poor wound healing, or numbness. Before starting high-dose antioxidant supplements, ask your doctor or pharmacist whether they interact with your medicines -- vitamin E, for example, can increase the effect of anticoagulants such as warfarin.

\textbf{Call 000 (or local emergency services) for:}
\begin{itemize}
    \item Severe or uncontrolled bleeding
    \item Difficulty breathing, facial swelling, or collapse after taking a supplement (allergic reaction)
    \item Chest pain, sudden confusion, or other symptoms that suggest a serious medical problem rather than a nutritional issue
\end{itemize}

\section*{Diagnosis and Investigations}
\begin{itemize}
    \item \textbf{Dietary history:} a review of food intake, including restrictive, elimination, or fad diets
    \item \textbf{Blood tests:} specific vitamin or mineral levels (such as vitamin C, vitamin E, selenium, or zinc) only when deficiency is clinically suspected
    \item \textbf{Full blood count and iron studies:} to detect anaemia that often accompanies or mimics nutritional problems
    \item \textbf{Medication and supplement review:} a complete list of prescribed medicines, over-the-counter products, and supplements to identify interactions
    \item \textbf{No routine testing:} measuring ``oxidative stress'' or antioxidant status in healthy people is not part of standard practice and is not clinically useful
\end{itemize}

\section*{Management}
\begin{itemize}
    \item \textbf{Food first:} a varied diet rich in colourful fruit and vegetables, nuts, seeds, and wholegrains supplies antioxidants naturally, alongside fibre and other nutrients
    \item \textbf{Supplements for specific deficiencies:} only when a genuine deficiency is diagnosed, such as vitamin C replacement for scurvy or zinc for confirmed deficiency
    \item \textbf{Evidence-based formulas for specific conditions:} for example, the AREDS2 formulation for age-related macular degeneration, but only under medical guidance
    \item \textbf{Avoid mega-doses:} high-dose supplements have not been shown to prevent major disease and can cause harm, including bleeding and toxicity
    \item \textbf{Address the drivers of oxidative stress:} stopping smoking, limiting alcohol, and protecting the skin from sun are more powerful than any supplement
\end{itemize}

\section*{Complications}
\begin{itemize}
    \item High-dose beta-carotene supplements increased the risk of lung cancer in smokers in clinical trials
    \item High-dose vitamin E has been associated with increased bleeding risk, including haemorrhagic stroke
    \item Selenium excess can cause toxicity with brittle hair and nails, and gastrointestinal upset
    \item Supplements can interact with prescribed medicines, notably anticoagulants and chemotherapy
    \item Relying on supplements may delay the diagnosis and treatment of the true cause of symptoms
    \item The antioxidant paradox: in excess, some ``antioxidants'' can act as pro-oxidants and increase cell damage
\end{itemize}

\section*{Prognosis and Outcomes}
For healthy people eating a varied diet, there is no good evidence that antioxidant supplements prevent heart disease, cancer, or dementia, or that they extend life. The consistent finding across large trials is that the protection associated with antioxidants in food is not reproduced by pills, and some high-dose supplements have increased, rather than reduced, risk.

When a true deficiency is identified and corrected, the outlook is excellent. The practical message is unchanged: a balanced diet rich in plant foods is the safest and most reliable way to obtain antioxidants, while the greatest health gains come from quitting smoking, limiting alcohol, maintaining activity, and controlling weight.

\section*{Prevention and Self-Care}
\begin{itemize}
    \item Eat a variety of colourful fruit and vegetables, nuts, seeds, legumes, and wholegrains each day
    \item Get vitamins from food rather than relying on high-dose pills
    \item Limit smoking and excess alcohol, which increase oxidative stress
    \item Protect skin from excess sun with clothing, shade, and sunscreen
    \item Keep up regular physical activity and a healthy weight
    \item If you choose supplements, tell your doctor and pharmacist, and stay within recommended doses
\end{itemize}

\section*{Sources and Further Reading}
The following organisations provide reliable, up-to-date information on antioxidants, nutrition, and supplements for patients, students, and health professionals.

\begin{itemize}
    \item NHS. \textit{Information on vitamins, minerals, and healthy eating.} https://www.nhs.uk/conditions/
    \item Mayo Clinic. \textit{Overview of antioxidants, supplements, and diet.} https://www.mayoclinic.org/
    \item MedlinePlus. \textit{Health information on vitamins, minerals, and dietary supplements.} https://medlineplus.gov/
    \item World Health Organization. \textit{Resources on nutrition and healthy diets.} https://www.who.int/
    \item MSD Manual. \textit{Professional information on vitamins, minerals, and nutritional disorders.} https://www.msdmanuals.com/
    \item National Institutes of Health Office of Dietary Supplements. \textit{Fact sheets on vitamins and supplements.} https://ods.od.nih.gov/
\end{itemize}
